%!TEX program = uplatex
\RequirePackage{plautopatch}
\documentclass[uplatex,dvipdfmx]{jlreq}
\usepackage{amsmath,amssymb}
\usepackage{enumerate}

\pagestyle{empty}

\begin{document}

%======================================================================
% 表紙
%======================================================================
\begin{center}
  \vspace*{20pt}
  {\Huge\bfseries ENCOUNTER with MATHEMATICS}\\[10pt]
  {\Large 第42回}\\[6pt]
  {\large 2007年11月30日、12月1日}\\[20pt]
  {\Huge\bfseries Euler 生誕300年}\\[6pt]
  {\Large Euler からゼータの世界へ}\\[30pt]
  {\bfseries 梗\qquad 概}
\end{center}

オイラー（1707年4月15日〜1783年9月18日）の生誕300年を記念して、オイラーの創始したゼータ研究の過去・現在・未来を考えたい。オイラーはゼータについて画期的な研究を行ない、その発展として、ゼータは現代数学の重要な研究分野となっている。

オイラーはスイスのバーゼルに生まれたが、20歳になった1727年にロシアのサンクトペテルブルグに移って、そこで生涯の大部分を送り、亡くなった。この集会では、オイラーのゼータ関係の仕事の概略とサンクトペテルブルグで行われたオイラーフェスティバルの様子を黒川が紹介する。

オイラーはゼータの収束値とともにゼータの発散値も求めることで、ゼータの関数等式を導いた。その際に、有名な微分作用素であるオイラー作用素を用いた。同時に、負の偶数が零点になることも発見した。この、関数等式と実零点の計算は「虚の零点はすべて実部が $1/2$ の対称軸上に乗るだろう」という数学最高の予想「リーマン予想」の前触れとなった。オイラー作用素からの発展を落合が解説する。

オイラーのゼータは、もともとは自然数全体に関する和から出発したが、素数全体に関する積にもなるというオイラーの大発見（オイラー積の発見）によって素数の研究に結びついた。リーマンがリーマン予想の提出に至ったのもゼータから素数分布を研究するためであった。現在までに、オイラー積は多方面に一般化されて研究されてきていて、数論の研究はオイラー積の研究と言える。その未来も含め平野が解説する。

オイラーはゼータの相棒である「保型形式」の創始者でもある。とくに、オイラーの五角数定理はイータ関数（無限積）とテータ関数（無限和）の一致を示しているもので、保型形式論の端緒と言える。保型形式は積分変換によってゼータを作り出す。これまでに、保型形式は群論や対称空間論や表現論に結びついて広大なゼータの世界を開拓してきていて、重要な未解決の問題も数多い。その様子を権が解説する。

\vfill

\newpage

%======================================================================
% 講演1：黒川信重「サンクトペテルブルクのオイラー」
%======================================================================
\begin{center}
{\large\bfseries サンクトペテルブルクのオイラー}\\[6pt]
\hspace*{\fill}{\large\bfseries 黒川　信重（東京工業大学・理工学研究科）}
\end{center}

オイラーは20歳で水の都サンクトペテルブルグに到着するや、10年のうちに、ガンマ関数やゼータ関数の研究に大成功を収めた。ちなみに「ガンマ」も「ゼータ」の仲間である。サンクトペテルブルグにはオイラーの数学研究の成功の秘密があったのだろうか。いずれにしても、オイラーにサンクトペテルブルグの水が合ったことは間違いない。

講演では、ガンマやゼータ等のサンクトペテルブルグにおけるオイラーの発見を中心に解説し、本年6月に白夜のサンクトペテルブルグで行われたオイラーフェスティバルの様子も紹介したい。
% 単純リー環にはベキ零元というやっかいな代物が住んでいます。ベキ零元がなければ、その構造ももう少しあつかいやすものになっていたでしょうが、存在するために複雑になりました。逆に言えば、そのために構造が豊富になりました。

\newpage

%======================================================================
% 講演2：落合啓之「オイラー作用素からの発展」
%======================================================================
\begin{center}
{\large\bfseries オイラー作用素からの発展}\\[6pt]
\hspace*{\fill}{\large\bfseries 落合　啓之（名古屋大学・多元数理科学研究科）}
\end{center}

オイラーはゼータの発散値、すなわち、負の整数点での値を求めている。最初のいくつかの例を挙げると、
\begin{align*}
\zeta(-1) &= 1+2+3+4+5+6+\cdots = -\frac{1}{12}, \\
\zeta(-3) &= 1+2^3+3^3+4^3+5^3+6^3+\cdots = \frac{1}{120}
\end{align*}
であり、また、負の偶数点での値は零である。これを示す際に、微分作用素 $x \frac{d}{dx}$ が用いられている。これがのちに有名になる\textbf{オイラー作用素}である。講演ではオイラーの計算を再現することから始め、オイラー作用素の行列式によるガンマ関数の表示、多変数化などに触れてみたい。また、リーマン予想の一つのアプローチである Meyer の仕事にも簡単に触れたい。

\begin{center}\textsc{参考文献}\end{center}
\begin{enumerate}[{[1]}]
  \item 黒川信重：『オイラー探検：無限大の滝と12連峰』シュプリンガージャパン、2007年9月.
  \item R. Meyer, A spectral interpretation for the zeros of the Riemann zeta function, arXiv:math.NT/0412277.
\end{enumerate}

\newpage

%======================================================================
% 講演3：平野幹「オイラー積の300年」
%======================================================================
\begin{center}
{\large\bfseries オイラー積の300年}\\[6pt]
\hspace*{\fill}{\large\bfseries 平野　幹（成蹊大学・理工学部）}
\end{center}

オイラーは1737年、ゼータ関数の素数全体に関する無限積表示
\[
\sum_{n=1}^\infty \frac{1}{n^s}=\prod_p\left(1-\frac{1}{p^s}\right)^{-1}
\]
を発見した。この表示は、$s=1$ で左辺の無限級数が発散することから、素数が無限に存在することの証明をもたらす。以来、この右辺の無限積はオイラー積と呼ばれ、素数論、特に素数分布論の基礎をなす最も重要なもののひとつとして位置づけられる。

また、オイラー積の一般化によって、素元が定義されるような様々な数論的対象にゼータ関数が定義され、その研究の歴史・成果は膨大である。

\begin{center}\textsc{参考文献}\end{center}
\begin{enumerate}[{[1]}]
  \item 第9回整数論サマースクール報告集「ゼータ関数」（2001年7月15日〜7月19日、於 休暇村大久野島）.
\end{enumerate}

\newpage

%======================================================================
% 講演4：権寧魯「オイラー五角数定理から保型形式へ」
%======================================================================
\begin{center}
{\large\bfseries オイラー五角数定理から保型形式へ}\\[6pt]
\hspace*{\fill}{\large\bfseries 権　寧魯（九州大学・数理学研究院）}
\end{center}

オイラーは1741年に五角数定理を発見した。五角数定理とは、無限積＝無限和の形をした等式で、右辺のべき級数の指数に「五角数」が現れることからこの名がついている。この定理はオイラーによって1750年に証明され、分割数、約数関数、発散級数などに関する研究に用いられた。その後、この定理はさまざまな形で一般化された（ヤコビの三重積公式、ロジャース・ラマヌジャン恒等式、ザギエ公式、……）。また、五角数定理の左辺の無限積＝デデキントのイータ関数であり、右辺の無限和＝ヤコビのテータ関数なので、五角数定理は「保型形式」たちの間に成り立つ等式と見なせる。その後、ラマヌジャンはラマヌジャンのデルタ関数（＝イータ関数の24乗）を無限和に展開し、その係数（＝ラマヌジャンのタウ関数）を調べることによって、2次のオイラー積を持つゼータ関数を発見した。このゼータ関数はラマヌジャンのデルタ関数（＝重さ12のカスプ形式）を用いて積分表示される。デルタ関数の保型性を用いて、このゼータ関数は複素数平面全体に正則に解析接続されて、ある関数等式をみたすことが証明できる。このように、五角数定理は保型形式を用いてエル関数やゼータ関数を研究するといった手法の糸口を与えたともいえる。

この講演の前半では、オイラーの五角数定理のオイラーによる証明を紹介し、そこから得られる分割関数や約数関数についてのいくつかの応用について述べる。後半では、保型形式の基礎的な理論について解説し、そこから得られるゼータ関数たちを紹介する。時間が許せば、五角数定理のいくつかの一般化や関連する話題についても触れる予定である。

\begin{center}\textsc{参考文献}\end{center}
\begin{enumerate}[{[1]}]
  \item G. Andrews, The theory of Partitions, Encyclopedia of Mathematics and its Applications, vol.~2, Addison-Wesley, 1976.
  \item J. Bell, Euler and the pentagonal number theorem, arXiv:math.HO/0510054v2.
  \item L. Euler, Evolutio producti infiniti $(1-x)(1-xx)(1-x^3)(1-x^4)(1-x^5)(1-x^6)$ etc.\ in seriem simplicem, Opera Omnia Series 1: Opera mathematica vol.~3.
  \item 梅田亨：跡等式としての五角数定理、数理解析研究所 講究録、vol.~1497、88--102.
  \item D. Zagier, Vassiliev invariants and a strange identity related to the Dedekind eta-function, Topology \textbf{40} (2001), 945--960.
\end{enumerate}

\end{document}
